\documentclass[17pt]{article}

\usepackage{cmap}	% поиск по документу
\usepackage{hyperref}	% поддержка гиперссылок
\usepackage{geometry}	% настройка геометрии документа
\usepackage{multicol}   % настройка шаблона расположения текста
\usepackage{extsizes}	% расширенная настройка размеров шрифтов
\usepackage{indentfirst}    % отступ для первого абзаца
\usepackage{tabularx}	% расширенные таблицы
\usepackage{makecell}   % настройка клеток таблиц
\usepackage{caption}    % настройка подписей к вставкам
\usepackage[table]{xcolor}  % покраска таблиц
\usepackage{graphicx}	% вставка изображений
\usepackage{wrapfig, booktabs}	% оформление таблиц
\usepackage{graphicx, caption, subcaption}
\usepackage{minted}    % вставка кода

\usepackage{color}	% настройка цвета текста - \textcolor
\usepackage{soul}	% выделение текста - \hl
\usepackage[T2A]{fontenc}	% поддержка кириллических символов
\usepackage[utf8]{inputenc} % поддержка кириллических символов
\usepackage[english,russian]{babel}	% поддержка кириллических символов
\usepackage{csquotes, biblatex}   % вывод библиографии

\usepackage{mathcmd}	% математические команды
\usepackage{amsfonts}	% математические символы 1
\usepackage{amssymb}	% математические символы 2
\usepackage{amsmath}	% математические символы 3
\usepackage{amsthm}	% работа с теоремами

% Настройка окружения теорем 1 (RU)
\newtheorem{theorem}{Теорема}[section]
\newtheorem{corollary}{Следствие}[theorem]
\newtheorem{lemma}[theorem]{Лемма}
\newtheorem{example}{Пример}[section]
\newtheorem*{remark}{Утверждение}
\addto\captionsrussian{\renewcommand\proofname{Док-во:}}

% Настройка окружения теорем 2 (EN)
%\newtheorem{theorem}{TH}[section]
%\newtheorem{corollary}{CO}[theorem]
%\newtheorem{lemma}[theorem]{LM}
%\newtheorem{example}{EX}[section]
%\addto\captionsrussian{\renewcommand\proofname{PF:}}

\renewcommand{\labelitemii}{$\circ$}
\renewcommand{\labelitemiii}{-}

\addbibresource{biblio.bib}
\graphicspath{ {./images/} }

% Настройка оформления вставок кода
\setminted{xleftmargin=\parindent, fontsize=\footnotesize, linenos, tabsize=2, breaklines}
\usemintedstyle{emacs} 

% Настройка отступов
\geometry{
    a4paper,
    % total={200mm,287mm},
    left=10mm,
    right=10mm,
    top=10mm,
    bottom=20mm
}

% Заголовок
\title{
    \textbf{Методы моделирования и анализа стохастических систем} \\
    \large Лабораторная работа № 5: \\
    Исследование индуцированных шумом переходов в модели SN
}
\author{Костарев Иван (МЕНМ-150201)}
\date{весна 2025/2026}


\begin{document}

\maketitle
\tableofcontents



\newpage
\section{Аттракторы детерминированной модели}
\label{sec:attractors}

(A) Изобразим 2 устойчивых равновесия при значениях параметра $a = 0.5$, $a = 7$ (Рис. \ref{fig:attractors_a}).

\begin{figure}[htbp]
    \centering
    \includegraphics[width=1\textwidth]{attractors_a.png}
    \caption{Аттракторы при $a = 0.5$, $a = 7$.}
    \label{fig:attractors_a}
\end{figure}

(B) Изобразим 2 устойчивых равновесия, а также устойчивый и неустойчивый предельный цикл при значениях параметра $a = 0.715$, $a = 0.75$ (Рис. \ref{fig:attractors_b}).

\begin{figure}[htbp]
    \centering
    \includegraphics[width=1\textwidth]{attractors_b.png}
    \caption{Аттракторы при $a = 0.715$, $a = 0.75$.}
    \label{fig:attractors_b}
\end{figure}

(C)$^*$ Изобразим 2 устойчивых равновесия, а также устойчивый и 2 неустойчивых предельных цикла при значениях параметра $a = 0.8$, $a = 0.9$ (Рис. \ref{fig:attractors_c}).

\begin{figure}[htbp]
    \centering
    \includegraphics[width=1\textwidth]{attractors_c.png}
    \caption{Аттракторы при $a = 0.8$, $a = 0.9$.}
    \label{fig:attractors_c}
\end{figure}



\newpage
\section{Исследование ИШП между устойчивыми равновесиями для случая A}

Исследуем и визуализируем переходы между устойчивыми равновесиями системы SN для случая A из раздела \ref{sec:attractors}, при $a = 0.5$.

(a) Локализуем бассейны притяжения равновесий $M_1$, $M_2$ (Рис. \ref{fig:basins_of_attraction}).

\begin{figure}[htbp]
    \centering
    \includegraphics[width=0.6\textwidth]{basins_of_attraction.png}
    \caption{Бассейны притяжения равновесий.}
    \label{fig:basins_of_attraction}
\end{figure}

(б) Построим доверительные эллипсы для равновесий $M_1$, $M_2$ (Рис. \ref{fig:confidence_ellipse}).
С ростом $\varepsilon$ доверительные эллипсы выходят за границы бассейна притяжения своего равновесия.

\begin{figure}[htbp]
    \centering
    \includegraphics[width=0.6\textwidth]{confidence_ellipse.png}
    \caption{Доверительные эллипсы для аттракторов при $\varepsilon = 0.1$, $\varepsilon = 0.2$.}
    \label{fig:confidence_ellipse}
\end{figure}

(в) Построим временные ряды для траекторий с началом в равновесиях $M_1$, $M_2$ (Рис. \ref{fig:time_series}).
С ростом $\varepsilon$, после достижения определенного порога, можно наблюдать индуцированные шумами переходы между равновесиями системы.

\begin{figure}[htbp]
    \centering
    \includegraphics[width=1\textwidth]{time_series.png}
    \caption{Временные ряды для траекторий при $\varepsilon = 0.1$, $\varepsilon = 0.2$.}
    \label{fig:time_series}
\end{figure}



\newpage
\section{Исследование ИШП между равновесиями и предельным циклом для случая C}

Исследуем и визуализируем переходы между устойчивыми равновесиями системы SN для случая C из раздела \ref{sec:attractors}, при $a = 0.9$.

(а) Построим доверительные эллипсы для равновесий $M_1$, $M_2$ (Рис. \ref{fig:cycle_confidence_ellipse}).
С ростом $\varepsilon$ эллипсы выходят за границы неустойчивого предельного цикла вокруг равновесия, что делает возможным выход траектории из окрестности равновесия на предельный цикл.

\begin{figure}[htbp]
    \centering
    \includegraphics[width=0.6\textwidth]{cycle_confidence_ellipse.png}
    \caption{Доверительные эллипсы для аттракторов при $\varepsilon = 0.02$, $\varepsilon = 0.05$.}
    \label{fig:cycle_confidence_ellipse}
\end{figure}

(б) Построим временные ряды для траекторий с началом в равновесиях $M_1$, $M_2$ (Рис. \ref{fig:cycle_time_series}).
С ростом $\varepsilon$, после достижения определенного порога, можно наблюдать индуцированные шумами переходы между равновесиями и неустойчивыми предельными циклами.

\begin{figure}[htbp]
    \centering
    \includegraphics[width=1\textwidth]{cycle_time_series.png}
    \caption{Временные ряды для траекторий при $\varepsilon = 0.02$, $\varepsilon = 0.05$.}
    \label{fig:cycle_time_series}
\end{figure}

\end{document}